\documentclass{ltjsbook}
\usepackage{docmute} 
\usepackage{../../myStyle} 

\begin{document}
\chapter{距離を分析する}
\label{H08_Distance}

ここまで，心理尺度の理論と実践を見てきました。理論的仮定から遠く離れて，実践的に「聞けばわかる」といわんばかりにさまざまな``こころ''の状態が測定されていることが明らかになってきました。

心がどのような状態なのか，言葉にすることが可能なのか，言葉にしたものは心とは違う何かなのか，といった哲学的な問題については，これから心理学業界全体を上げて考え直さなければならない問題なのかもしれません。あるいはこれまでのように，理論的正当性は欠いたままかもしれないが，測られたものによって言論空間が成立し，何かしら妥当性を感じさせる議論ができればそれでよい，というプラグマティックな考え方を選択することになるかもしれません。

こうした問題について，この講義の中で答えが出せるわけではありませんが，方法論的な改善という観点からアプローチしてみたいと思います。
心理尺度の使われ方の問題は，等間隔性を仮定して数値化されたデータに対し，因子分析によって潜在変数を\kenten{見出し}，見出された構成概念が実在するものと考えて論じるという一連の研究実践で醸造されてくるものでした。
等間隔性の仮定や因子分析という手法については，測定モデルを工夫することで修正できる可能性が残されています。

以下では測定されたものが順序尺度水準や名義尺度水準といった，より低い尺度水準の値であることを想定したモデルの応用を考えていきたいと思います。実は心理学において，因子分析による尺度作成という慣例ができあがる前は，さまざまな分析手法が提案，適用されてきました。今一度これらの分析モデルに目を向けることで，新しい可能性について考えてみたいと思います。

\section{距離と心理学のデータ}
\label{distance}
私たちが単位もない不確かなものを対象にしながらもそれを測定するモノサシをつくることができるのは，2つの対象にたいしてその比較をして一方が他方よりも大である，ということが言えるからでしょう。
記号を使って言えば，$x$と$y$を比べて$x \succeq y$である，ということから，尺度上の値$p \ge q$を対応させるということが，尺度を作るということです\footnote{経済学では学問の最初の段階で，財を源とした集合に対する二項関係$x \circ y$を定義し，それを効用関数$u$で実数領域に写像して$u(x) \ge u(y)$を考える，といった基本的な原理を抑えます。心理学ではなぜか，「集合」「元」「二項関係」という基本的な比較の要素について考えることなく，その分析ツールだけがどんどんと発展してきています。}。
このとき比較する2つが重さや広さのような物理的なものであればわかりやすいですし，「痛み」や「喜び」といった心理的な要素であっても構いません。あるいは「より賛成」といった態度表明のようなものでも良いかもしれません。
この比較から出てくる関係は，分散共分散や相関係数のように強い線形の仮定を置かなくても，より緩やかに\keyterm{距離}{きょり}{distance}という考え方で表現されます\footnote{後述しますが，相関係数も距離の一種と考えられます。}。

あらためて\keyterm{距離}{きょり}{distance}とは何かを考えてみましょう。2点$x,y$の距離を$d(x,y)$とすると，距離とよばれる数字の条件は次のようになります。
\begin{description}
    \item[非負性] $d(x,y)\ge 0$
    \item[非退化性] $d(x,y)=0 \Leftrightarrow x=y$
    \item[対称性] $d(x,y)=d(y,x)$
    \item[三角不等式] $d(x,z) + d(z,y) \ge d(x,y)$
\end{description}
もっとも一般的に使われるのはユークリッド距離で，2次元座標$(x_1,y_1),(x_2,y_2)$があった時のユークリッド距離は次のように計算します。
\[ \sqrt{(x_1-x_2)^2+(y_1-y_2)^2} \]
3次元座標$(x_1,y_1,z_1),(x_2,y_2,z_2)$の場合は項を増やすだけです。
\[ \sqrt{(x_1-x_2)^2+(y_1-y_2)^2+(z_1-z_2)^2} \]
このようにして計算される距離ですが，上の条件を考えると何も二乗してルートを取らなくても絶対値を足し合わせるような方法でも構いません。たとえば2次元座標の場合は，次のような計算でもいいのです。
\[ |x_1-x_2|+|y_1-y_2| \]
現にこのような距離のことを\keyterm{マンハッタン距離}{まんはったんきょり}{Manhattan distance}といいます。二乗ではなくn乗してn乗根をとる，という形で一般化することもできます。
\[ d(x,y) = \left( \sum_{i=1}^n |x_i - x_j|^p\right)^{\frac{1}{p}} \]
これは\keyterm{ミンコフスキー距離}{みんこふすきーきょり}{Minkowski distance}という名前がついています。他にもマキシマム距離，バイナリ距離，チェビシェフの距離，キャンベラ距離などがあり，いずれもRの\texttt{dist}関数のオプションで選ぶことができますので，ヘルプなどを参照してみてください。

また，相関係数$r_{ij}$は$-1 \le r_{ij} \le +1$の範囲にありますが，この絶対値から$1-|r_{ij}|$とするとこれも距離の条件に当てはまります。どの程度類似しているかというのも，距離と考えることができるのです。

ここまでは数学的な距離のバリエーションでしたが，次に心理学的な意味に目を向けてみましょう。
距離とは類似度でもありますから，何を距離と見なすかによって，さまざまな心理学的刺激がデータを形作ることになります。

たとえば評定尺度で，n個の項目で何らかの評価をしてもらったとします。$x_{ij}$を$i$番目の項目における対象$j$の評定値だとすると，
$d(j,k) = \sqrt{\sum_{i=1}^N (x_{ij}-x_{ik})^2 }$とすれば対象の類似度が計算できます。


(非)類似性を距離とみなして用いることが一般的
(非)類似性の心理データを取る利点
被験者の自然な(総合的な)判断に任せられる；下位の評価次元を実験者が準備することは研究者があらかじめどの属性が重要かを決め打ちしているようなもの
上と関連して，人の評価傾向(社会的望ましさバイアスなど）から自由なデータを入手できる
被験者には何らかの内的一貫性を求めており，データ収集時にその一貫性をチェックすることができる

尺度評定を用いる方法
刺激の混同率
代替価/連想価
刺激の汎化勾配Si→Riの条件づけがあるとき，Pr(Ri|Sj)を類似性とする
反応潜時２つの刺激が「同じ」か「違う」かを判断する課題への反応潜時を類似性の指標とする
ソシオメトリックなデータ



ほかにも何らかの刺激A,Bについて，AかBかの判断をさせた時に混同してしまった混同率や，単語と単語の連想価，刺激の汎化勾配，反応潜時，ソシオメトリックなデータなど，いろいろなものが「類似しているかどうか」の指標として使えます\footnote{これらデータの例に関しては\citet{Takane198002}のPp.14--27を参照してください。}。尺度評定よりも，具体的な2つの刺激が似ているかどうかの反応の方がやりやすい，というのは誰しも実感としてわかることではないかと思います。

類似度のデータが得られれば，それを距離と見做してMDSにかければ，変数間の関係を地図に描くことができるわけです。このように応用可能な領域が非常に広いことも，MDSの利点であると言えるでしょう。

\section{非計量多次元尺度法}

さて心理学的なさまざまな刺激が，MDSによって可視化できるということがわかってきました。距離行列が構成できれば，後の分析は何とでもできるわけです。
ただし，この場合の距離行列とは，間隔尺度水準以上であることが必要です。
しかし心理学的な刺激に対する反応をデータ化する時は，同時に被験者はそこまで鋭敏に反応しているのか，いいかえれば本当に間隔尺度水準以上の精度で判断できているのかな，という人間側の問題が気になりますね。そこまで人間は鋭敏無反応をしていないかもしれません。

でも大丈夫。MDSは仮定を緩めたモデルがあります。\keyterm{非計量的多次元尺度構成法}{ひけいりょうてきたじげんしゃくどこうせいほう}{Non-Metric Multi-Dimensional Scaling}と呼ばれる手法がそれです。非計量MDSでは，対象$j$と$k$の類似度を$\delta_{jk}$とし，分析によって埋め込む多次元空間での距離を$d_{jk}$とすると，

$$ \delta_{jk} > \delta_{lm} \mbox{ならば}d_{jk}\le d_{lm} $$

のように，イコールではなく順序関係だけ保持して座標を求めます。この手法では，元データが順序尺度水準程度の情報しか持っていなくても地図を描くことができるのです。人間の判断はせいぜいが順序尺度ぐらいですから，「より類似している($\delta_{jk} > \delta_{lm}$)」のであれば「より近くにある($d_{jk}\le d_{lm} $)」というぐらいの配置の方が良いかもしれません。


表\ref{tbl::15_01}に示したような物理的距離であれば，間隔尺度水準の情報であることは間違いありません。その場合には，最初に示した固有値分解による手法を使います。これは非計量MDSに対して，\keyterm{計量的多次元尺度構成法}{けいりょうてきたじげんしゃくどこうせいほう}{Metric Multi-Dimensional Scaling}と呼ばれています。

\section{非計量多次元尺度法}

さて心理学的なさまざまな刺激が，MDSによって可視化できるということがわかってきました。距離行列が構成できれば，後の分析は何とでもできるわけです。
ただし，この場合の距離行列とは，間隔尺度水準以上であることが必要です。
しかし心理学的な刺激に対する反応をデータ化する時は，同時に被験者はそこまで鋭敏に反応しているのか，いいかえれば本当に間隔尺度水準以上の精度で判断できているのかな，という人間側の問題が気になりますね。そこまで人間は鋭敏無反応をしていないかもしれません。

でも大丈夫。MDSは仮定を緩めたモデルがあります。\keyterm{非計量的多次元尺度構成法}{ひけいりょうてきたじげんしゃくどこうせいほう}{Non-Metric Multi-Dimensional Scaling}と呼ばれる手法がそれです。非計量MDSでは，対象$j$と$k$の類似度を$\delta_{jk}$とし，分析によって埋め込む多次元空間での距離を$d_{jk}$とすると，

$$ \delta_{jk} > \delta_{lm} \mbox{ならば}d_{jk}\le d_{lm} $$

のように，イコールではなく順序関係だけ保持して座標を求めます。この手法では，元データが順序尺度水準程度の情報しか持っていなくても地図を描くことができるのです。人間の判断はせいぜいが順序尺度ぐらいですから，「より類似している($\delta_{jk} > \delta_{lm}$)」のであれば「より近くにある($d_{jk}\le d_{lm} $)」というぐらいの配置の方が良いかもしれません。


表\ref{tbl::15_01}に示したような物理的距離であれば，間隔尺度水準の情報であることは間違いありません。その場合には，最初に示した固有値分解による手法を使います。これは非計量MDSに対して，\keyterm{計量的多次元尺度構成法}{けいりょうてきたじげんしゃくどこうせいほう}{Metric Multi-Dimensional Scaling}と呼ばれています。
\end{document}